To understand the design constraints and operational mechanics of the proposed translation pipeline, we first review the architectural constraints of the target BCI hardware, the computational paradigms introduced by RVV, and the foundational principles of the compilers used in this research. 

The fundamental engineering challenges in the development of implantable BCIs is the strict management of thermal dissipation. The human cerebral cortex is highly sensitive to temperature fluctuation, and prolonged exposure to high temperatures can cause tissue damage. Consequently, international medical standards dictate that an implant must not raise the local tissue temperature by more than one degree Celsius, effectively restricting the total power consumption of the entire BCI system—including amplifiers, analog-to-digital converters (ADCs), wireless telemetry, and digital signal processing (DSP) units—to roughly fifteen milliwatts \cite{10120914}.

Historically, BCI designers have relied on custom-designed ASICs to meet these power requirements \cite{7982670} \cite{Musk2019} \cite{Heck2014}. While ASICs offer high energy efficiency, they inherently lack flexibility. If a neuroscientist develops a superior seizure-prediction algorithm or wishes to repurpose the implant for a different disorder, a completely new physical chip must be engineered, fabricated, and implanted \cite{9138938}. General purpose microprocessors offer the required flexibility through software updates but consume power orders of magnitude beyond the safe physiological limit \cite{10120914}.

HALO resolves this through aggressive hardware-software co-design. The architecture is built around a configurable interconnect that links a heterogeneous array of application specific PEs \cite{9138938}. These PEs are optimized for the specific computational  kernels most frequently encountered in neural signal processing, such as FFT, XCOR, and Butterworth bandpass filtering \cite{10120914}. These accelerators can be dynamically chained together to form highly efficient data flow pipelines tailored to specific clinical tasks \cite{9138938}.

To manage these pipelines and execute computational tasks that fall outside the capabilities of a standalone PE, HALO uses a low power RISC-V microcontroller \cite{9138938}. To ensure that this software-driven execution does not violate the power budget, the microcontroller relies heavily on RVV. Unlike traditional single instruction, multiple data (SIMD) architectures, which utilize fixed width vector registers and require code rewrites when the vector width changes for different hardware generations, RVV introduces a vector-length agnostic programming model \cite{pipicelli2020quantum, Ibanez}.

In the context of BCI workloads, which involve matrix multiplications and convolutions over various window sizes, RVV allows the processor to load large, contiguous segments of multi-channel neural data into vector registers and apply a single arithmetic instruction across the entire dataset simultaneously. This significantly reduces the dynamic power consumed by the processor’s instruction fetch and decode units. By replacing hundreds or thousands of scalar loop iterations with a single vector instruction, the processor minimizes the total number of active clock cycles required to complete a given task. Consequently, computational algorithms execute much faster, allowing the microcontroller to aggressively switch to low-power sleep states between neural data sampling intervals, which is crucial for extending the lifespan of the implantable battery. 

Despite these hardware features, the platform remains inaccessible to neuroscientists. Clinical algorithms are developed using high-level, dynamically typed languages like MATLAB, which abstract away memory management and hardware-specific optimizations that are crucial for running the algorithms on silicon. The automated pipeline proposed in this research bridges this gap by integrating two existing software repositories. The first is ts-traversal, a TypeScript-based source-to-source compiler that parses MATLAB syntax, infers static types, and generates functional C code \cite{ts}. The second is the RISCV-Matrix C library, which maps standard mathematical operations directly to RVV assembly instructions \cite{riscvmatrix}. Critical fixes have been made in both the translation logic and the vectorization coverage.


